\documentclass[11pt]{article}

\usepackage[margin=1in]{geometry}
\usepackage{amsmath,amssymb}
\usepackage{array}

\title{Three-Player Bilevel Generalized Nash Benchmark}
\author{}
\date{}

\begin{document}
\maketitle

\section*{Problem Statement}

This document writes the three-player bilevel benchmark used in the test suite.
There are three players, indexed by \(i \in \{1,2,3\}\). Player \(i\)'s
decision vector is
\[
  x_i = (x_{i,1},x_{i,2},x_{i,3},x_{i,4},x_{i,5}) \in \mathbb{R}^5,
\]
where \(x_{i,j}\) denotes the \(j\)-th decision variable of player \(i\). In the
code this is `x[Block(i)][j]`.

The bilevel preference order is:
\[
  \text{upper level controls coordinates } \{4,5\}, \qquad
  \text{lower level controls coordinates } \{1,2,3\}.
\]

\section*{Raw Objective Targets}

The raw objective target is the point each player's objective would choose if
there were no constraints. The constraints then force the final equilibrium away
from this raw target.

The known equilibrium is
\[
\begin{aligned}
  x_1^\star &= (0.40,\;0.90,\;0.55,\;1.20,\;0.80),\\
  x_2^\star &= (0.45,\;1.00,\;0.60,\;1.30,\;0.90),\\
  x_3^\star &= (0.50,\;1.10,\;0.65,\;1.40,\;1.00).
\end{aligned}
\]

The raw target for player \(i\) is
\[
  t_i = x_i^\star + (0.6,0,-0.4,0,0).
\]
Thus
\[
\begin{aligned}
  t_1 &= (1.00,\;0.90,\;0.15,\;1.20,\;0.80),\\
  t_2 &= (1.05,\;1.00,\;0.20,\;1.30,\;0.90),\\
  t_3 &= (1.10,\;1.10,\;0.25,\;1.40,\;1.00).
\end{aligned}
\]

\section*{Quadratic Bilevel Problem}

For each player \(i\), define the upper-level and lower-level objectives
\[
  J_i^{\mathrm{up}}(x_i)
  =
  \sum_{j \in \{4,5\}} (x_{i,j}-t_{i,j})^2,
\]
and
\[
  J_i^{\mathrm{low}}(y_i)
  =
  \sum_{j \in \{1,2,3\}} (y_{i,j}-t_{i,j})^2.
\]

The bilevel problem for player \(i\) is
\[
\begin{aligned}
  \min_{x_i \in \mathbb{R}^5}
  \quad & J_i^{\mathrm{up}}(x_i)\\
  \text{subject to}
  \quad &
  x_i \in
  \operatorname*{arg\,min}_{y_i \in \mathbb{R}^5}
  \left\{
    J_i^{\mathrm{low}}(y_i)
    \;:\;
    g_i(y_i,x_{-i}) \ge 0
  \right\},
\end{aligned}
\]
where \(x_{-i}\) denotes the decision vectors of the other two players.

\section*{Coupled Generalized Nash Constraints}

Let
\[
  \rho = 0.25
\]
and define the coupled resource usage
\[
  c_{i,j}(x)
  =
  x_{i,j} + \rho \sum_{k \ne i} x_{k,j}.
\]
Here \(c_{i,j}(x)\) is a function of the current players' decisions. The
uppercase constants below are fixed capacities or floors.

Each player has four coupled inequalities:
\[
  g_i(x) =
  \begin{bmatrix}
    C_{i,1} - c_{i,1}(x)\\
    C_{i,2} - c_{i,2}(x)\\
    c_{i,3}(x) - F_{i,3}\\
    c_{i,4}(x) - F_{i,4}
  \end{bmatrix}
  \ge 0.
\]

The constants are chosen from the known solution \(x^\star\):
\[
\begin{aligned}
  C_{i,1} &= c_{i,1}(x^\star),\\
  C_{i,2} &= c_{i,2}(x^\star) + 1,\\
  F_{i,3} &= c_{i,3}(x^\star),\\
  F_{i,4} &= c_{i,4}(x^\star) - 1.
\end{aligned}
\]

Numerically,
\[
\begin{array}{c|cccc}
  i & C_{i,1} & C_{i,2} & F_{i,3} & F_{i,4}\\
  \hline
  1 & 0.6375 & 2.4250 & 0.8625 & 0.8750\\
  2 & 0.6750 & 2.5000 & 0.9000 & 0.9500\\
  3 & 0.7125 & 2.5750 & 0.9375 & 1.0250
\end{array}
\]

\section*{Expected Solution and Active Set}

At the expected equilibrium \(x^\star\),
\[
  C_{i,1} - c_{i,1}(x^\star) = 0,
  \qquad
  c_{i,3}(x^\star) - F_{i,3} = 0.
\]
Thus the active constraints for each player are:
\[
  C_{i,1} - c_{i,1}(x) \ge 0,
  \qquad
  c_{i,3}(x) - F_{i,3} \ge 0.
\]

The inactive constraints have slack equal to \(1\):
\[
  C_{i,2} - c_{i,2}(x^\star) = 1,
  \qquad
  c_{i,4}(x^\star) - F_{i,4} = 1.
\]

The constraints are genuinely coupled because every \(c_{i,j}(x)\) depends on
all three players' decision variables. The active coordinate-1 and coordinate-3
resource equations are player-specific, so they pin the distribution across
players rather than only aggregate totals.

\section*{Nonlinear Variant Used in Tests}

The nonlinear test case keeps the same known solution and active set. It replaces
the squared-distance objective by
\[
  \sum_{j \in S} \left(\cosh(x_{i,j}-t_{i,j}) - 1\right),
\]
and replaces each linear inequality residual \(r(x) \ge 0\) by
\[
  \exp(r(x)) - 1 \ge 0.
\]
This preserves the feasible set and active set because
\[
  \exp(r)-1 = 0 \iff r=0,
  \qquad
  \exp(r)-1 > 0 \iff r>0.
\]

\end{document}
